%% =====================================================================
%%  Corrigés — Bilan 5 « Hier et aujourd'hui »
%% =====================================================================
\corrigesBilan{5}

\begin{corrige}{1}
\begin{reponses}[4]
\item \rep{was}
\item \rep{were}
\item \rep{Were}
\item \rep{wasn't}
\end{reponses}
\end{corrige}

\begin{corrige}{2}
\begin{reponses}[4]
\item \rep{watched}
\item \rep{studied}
\item \rep{went}
\item \rep{bought}
\end{reponses}
\end{corrige}

\begin{corrige}{3}
\begin{reponses}[1]
\item \rep{We didn't go to the beach.}
\item \rep{Did she call you?}
\item \rep{I didn't see the movie.}
\item \rep{Were they happy?}
\item \rep{He wasn't at home.}
\item \rep{Did you take the bus?}
\end{reponses}
\pourquoi{avec \en{be} (\en{was}, \en{were}), pas de \en{did} : on inverse ou on ajoute \en{not}. Avec
les autres verbes, \en{did} / \en{didn't} + base verbale.}
\end{corrige}

\begin{corrige}{4}
\begin{reponses}[2]
\item \rep{visited}
\item \rep{is reading}
\item \rep{have}
\item \rep{are going to fly}
\item \rep{made}
\item \rep{was}
\end{reponses}
\pourquoi{\en{Last year}, \en{two days ago}, \en{yesterday} : prétérit ; \en{Right now} : présent en
\en{-ing} ; \en{Every Sunday} : présent simple ; \en{Tomorrow} : \en{be going to}.}
\end{corrige}

\begin{corrige}{5}
\begin{reponses}[3][french]
\item \rep{V}
\item \rep{F}
\item \rep{V}
\item \rep{F}
\item \rep{F}
\item \rep{V}
\end{reponses}
\pourquoi{\en{flew} est le prétérit de \en{fly} (prendre l'avion). Ils ont rendu visite à leurs parents
et vu un match de baseball, que les Braves n'ont pas gagné. Texte de la piste
page~\pageref{trans-B5.01}.}
\end{corrige}

\begin{corrige}{6}
\begin{reponses}[3][french]
\item \rep{V}
\item \rep{F}
\item \rep{V}
\item \rep{F}
\item \rep{V}
\item \rep{F}
\end{reponses}
\pourquoi{le journal est daté du 16 mars et parle d'« hier » : l'anniversaire était donc le 15. La fête a
eu lieu au café de Sam, et c'est Sam qui a fait le gâteau.}
\end{corrige}

\begin{corrige}{7}
\begin{reponses}[1]
\item \rep{When were you born?}
\item \rep{Where did you go?}
\item \rep{What are you doing?}
\item \rep{What are you going to do?}
\end{reponses}
\end{corrige}

\begin{corrige}{8}
Réponse modèle :
\begin{anglais}
\begin{tcolorbox}[otBase,colback=otGreenTint,colframe=otGreenTint,fontupper=\fontsize{10.4}{14}\selectfont]
Last year, I visited Scotland with my family. It was beautiful, but it rained every day! Every
morning, I listen to English songs on the bus. Next year, I'm going to take an English class.
\end{tcolorbox}
\end{anglais}
\retiens{0,5 point par phrase correcte : un temps différent pour chaque idée, bien conjugué.}
\end{corrige}

\begin{corrige}{9}
Réponses modèles :
\begin{reponses}[1]
\item \rep{My name is Julie Moreau. I'm thirty-two and I live in Lille.}
\item \rep{I was born in Lille.}
\item \rep{Last year, I went to Italy with my husband.}
\item \rep{I read, I cook and I play tennis.}
\item \rep{I'm going to visit New York!}
\end{reponses}
\retiens{vous venez de terminer le niveau A1 : félicitations ! Le tome 2 vous attend pour aller plus loin.}
\end{corrige}
